\documentclass[manuscript,review,anonymous]{acmart}
\acmJournal{TEAC}
\setcopyright{none}
\usepackage{amssymb,mathtools,booktabs}
\IfFileExists{libertine.sty}{}{\microtypesetup{expansion=false}}
\newcommand{\PNE}{\mathrm{PNE}}
\newcommand{\OPT}{\mathrm{OPT}}
\newcommand{\F}{\mathcal F}
\newcommand{\ind}{\mathbf 1}
\newcommand{\xor}{\mathbin{\oplus}}
\newcommand{\A}{\mathsf A}
\newcommand{\B}{\mathsf B}
\newcommand{\R}{\mathbb R}
\newcommand{\N}{\mathbb N}
\DeclareMathOperator{\pos}{pos}
\title{Structural Methods for Approximate Stability in Weighted Congestion Games}
\ccsdesc[500]{Theory of computation~Algorithmic game theory and mechanism design}
\ccsdesc[300]{Theory of computation~Approximation algorithms analysis}
\begin{document}
\begin{abstract}
We study how local cost comparisons determine approximate stability in weighted congestion games with polynomial latencies. We develop tools for comparing endpoint costs with area estimates, combining resource potentials with a common scale for each player, and realizing stability on a structured set of states by legal congestion resources. These tools explain both the amplification caused by large atoms and the existence guarantees available under small atomic shares, compatible weight types, or bounded weight depth. As an application, we construct rational games of degree $d$ with no $(d-1200\lceil\log_2d\rceil^2)$-approximate pure Nash equilibrium, for every integer $d\ge2^{20}$. Together with the known $d$-approximate existence theorem, this answers the asymptotic existence question raised by Christodoulou, Gairing, Giannakopoulos, Po\c{c}as, and Waldmann. We also give conditions under which the same realization method preserves social optimization and approximation quality.
\end{abstract}
\keywords{weighted congestion games, approximate pure Nash equilibria, approximate potentials, social cost}
\maketitle

\section{Introduction}
Weighted congestion games describe systems in which indivisible users choose resources and the cost of each resource depends on its total load. Approximate pure Nash equilibria provide a finite-state notion of stability when exact equilibria may fail to exist. Their analysis often begins with a local comparison between a player's cost and a change in a resource potential. To obtain a statement about the whole game, however, these local comparisons must remain consistent across every resource that the player can enter or leave.

We study three aspects of this passage from local costs to global stability. First, we compare the cost at a polynomial's endpoint with a discrete approximation to the area beneath it. The comparison distinguishes the total mass in a window from the size of the individual atoms used to sample it. Second, we identify when local resource potentials admit the same multiplier for a player on all resources. For resources with two optional weight types, this compatibility is described by product inequalities around type cycles. Third, we give a construction that confines every stable state to a prescribed set of structured binary states, while retaining an explicit formula for the approximation tolerance introduced by the added resources.

Each tool has a positive use. Small actual atomic shares make the integral potential effective. On forests of interacting weight types, the common-scale construction gives an approximation factor of order $\sqrt d$, independently of the number of types. For arbitrary resource populations, a common ordering by decreasing weight gives an existence bound determined by the height of the player graph. The realization construction allows stable states to be selected on the smaller structured domain; with an additional externality bound, it also allows social optimization on that domain.

These comparisons also explain how large approximation factors arise. A short, finely sampled window can preserve an endpoint-to-area amplification of order $d$ as costs propagate between weight classes, even though the light users pay an area rather than an endpoint cost. We use this distinction to build a finite family in which every structured state has an improving deviation. The realization theorem then excludes all other states without increasing the degree.

\paragraph{Relation to the existence question.}
Caragiannis and Fanelli~\cite[Theorem~2]{CF2021} prove a $d$-approximate potential guarantee for positive-degree monomial latencies. Their marginal bounds extend it to nonnegative-coefficient polynomials of degree at most $d\ge1$: split each resource into one copy per positive term, so every strategy uses all its copies together, and give a constant term $a_{e,0}$ the potential $a_{e,0}$ times its load. This preserves all player costs and the common marginal interval from $w_i c_e(Y)/d$ to $w_i c_e(Y)$. Christodoulou, Gairing, Giannakopoulos, Po\c{c}as, and Waldmann~\cite[Section~7]{CGGPW2023} establish a nonexistence bound of order $\sqrt d/\log d$ and ask how to close the gap to the linear existence bound. Our construction gives a lower bound of $d-O((\log d)^2)$, and therefore determines the leading asymptotic constant of the universal threshold. The $d$ upper bound follows from the cited potential; the construction and the structural statements used here are proved below.

Caragiannis and Fanelli also obtain an $e^{1/\tau}$-approximate potential for their $\tau$-congested games with $\tau>0$~\cite[Corollary~5]{CF2021}: in their monomial model, at every state each user of a degree-$k_e$ resource shares it with other users of total weight at least $\tau k_ew_i$. They give an upper bound $(d+1)/(d+\delta)$ on the $(d+\delta)$-approximate price of stability for $0\le\delta\le1$~\cite[Theorem~7]{CF2021}. Our small-share bound gives an explicit integral factor under its own all-state share condition; our social-cost results select within a realized structured domain and require a separate externality condition to compare with the full-game optimum. Neither use of potential methods is claimed as new by itself.

The counterexamples and the positive results answer different questions. In the unrestricted class, the leading linear factor is unavoidable. Under verifiable conditions on actual shares, compatible scales, or decreasing weight depth, the same local comparisons yield smaller factors. A failed scale certificate concerns the specified potential representation; a full-state counterexample requires the separate realization argument.

\subsection{Model and equilibrium}\label{sec:model}
A finite weighted congestion game has players $I$, resources $E$, positive weights $w_i$, and nonempty finite strategy families $S_i\subseteq2^E$. At a state $s\in\prod_iS_i$ put
\[
 x_e(s)=\sum_{i:e\in s_i}w_i,\qquad
 C_i(s)=\sum_{e\in s_i}c_e(x_e(s)),\qquad
 c_e(x)=\sum_{k=0}^d a_{e,k}x^k,\quad a_{e,k}\ge0.
\]
A state is an $\alpha$-$\PNE$, $\alpha\ge1$, if $C_i(s)\le\alpha C_i(t_i,s_{-i})$ for all players and legal deviations. We use the margin $C_i(s)-\alpha C_i(t_i,s_{-i})$, including at zero costs. Multiplying a player's whole cost by its fixed positive weight leaves these stability comparisons unchanged. The social objective used below is specified separately:
\begin{equation}\label{eq:welfare}
 W(s)=\sum_iw_iC_i(s)=\sum_ex_e(s)c_e(x_e(s)).
\end{equation}
For a game $G$, write $\OPT(G)=\min_{s\in\prod_iS_i}W(s)$, with the minimum taken over all its pure states.
Unless stated otherwise, theorems allow real data; rational input and output are required for the explicit lower-bound games and exact finite algorithms.

A resource can have compulsory users: those whose every strategy contains it. Absorb their load $b_e$ into $\widehat c_e(z)=c_e(b_e+z)$. The function remains a nonnegative-coefficient polynomial of degree at most $d$. Its remaining potential users are those who can both use and omit the resource. A player may be compulsory on one resource and optional on another. Its compulsory resources are retained in every deviation, so this resourcewise preprocessing preserves all its comparisons. Zero-cost resources can be omitted from scale constraints.


\paragraph{Organization.}
Section~\ref{sec:tools} develops the endpoint estimates, common player scales, weight-depth bound, and realization theorem. Section~\ref{sec:application} combines them into explicit all-state counterexamples and deduces the asymptotic threshold. Section~\ref{sec:positive} applies the tools to small shares and to the selection of stable states with social-cost guarantees. Appendix~\ref{app:construction} supplies the parameter estimates and the complete phase analysis for the construction.

\paragraph{Conventions and notation.}
All games are finite, weights are positive, and latency polynomials have nonnegative coefficients. A degree bound is an upper bound unless we specify pure degree-$d$ monomials. All approximation factors are at least one, and strict improvement means a strict multiplicative cost reduction. We use the margin formulation also at zero cost. Logarithms are natural unless a base is shown. The finite constructions use rational data; their output size and improvement-path length are not asserted to be polynomial in a compact description.

\section{Structural tools for approximate stability}\label{sec:tools}
We first isolate the marginal condition that lets resource functions sum to a potential for the whole game.

\subsection{From resource marginals to a common player scale}
The approximate-potential argument of Caragiannis and Fanelli~\cite{CF2021} compares the insertion marginal of a resource function with the player's resource cost. For a deviation involving several resources, we need the comparison to use one multiplier for the moving player. The following lemma states this requirement explicitly.

\begin{lemma}[Common player scales]\label{lem:glue}
Suppose resource set functions $H_e$ and positive numbers $q_i$ satisfy the following bounds. For each resource $e$, optional player $i$, and local optional-user set occurring in a legal state where $i$ omits $e$, insert $i$ into that set and require
\begin{equation}\label{eq:mastermarginal}
 q_i\widehat c_e(Y)\le\Delta_iH_e\le\beta q_i\widehat c_e(Y),
\end{equation}
where $Y$ is the optional load after insertion. The same $q_i$ must work on all resources that player $i$ can toggle. Then $\Phi(s)=\sum_eH_e(s)$ satisfies, for every unilateral deviation $s\to t$ by $i$,
\begin{equation}\label{eq:masterdeviation}
 \Phi(t)-\Phi(s)
 \le q_i[\beta C_i(t)-C_i(s)-(\beta-1)C_{i,H}]
 \le q_i[\beta C_i(t)-C_i(s)],
\end{equation}
where $C_{i,H}$ is the common cost of the resources retained by the player. Every strict $\beta$ improvement decreases $\Phi$, every such improvement path terminates, and a $\beta$-$\PNE$ exists.
\end{lemma}
\begin{proof}
Fix a unilateral deviation by player $i$. We apply the lower marginal bound to every deleted resource in its old user set, and the upper marginal bound to every added resource in its new user set. Retained resources have unchanged load and potential. Summing these inequalities gives $q_i(\beta C_{i,\mathrm{add}}-C_{i,\mathrm{del}})$, which equals the first expression in~\eqref{eq:masterdeviation}. A strict $\beta$ improvement therefore strictly decreases $\Phi$. Since the state space is finite, every such improvement path terminates, and a global minimum of $\Phi$ is a $\beta$-$\PNE$.
\end{proof}
The insistence on a common $q_i$ is essential to this summation. All positive tools in the paper construct such scales; the lower-bound construction exploits a different propagation of endpoint and area costs.

\subsection{Endpoint costs and area estimates}\label{sec:area}
The following estimate separates the mass in a window from its sampling resolution. It will control the cost transmitted between adjacent weight classes in Section~\ref{sec:application}.
Let $0=x_0<\cdots<x_n=Q$, $a_j=x_j-x_{j-1}$, $T=b+Q$, and $h=\max_ja_j$, with $b\ge0$ and $Q>0$.
\begin{lemma}[Endpoint--area error]\label{lem:area}
For every nondecreasing function $f$ on $[b,T]$,
\begin{equation}\label{eq:areaerror}
 \int_b^Tf(t)\,dt\le\sum_ja_jf(b+x_j)
 \le\int_b^Tf(t)\,dt+h[f(T)-f(b)].
\end{equation}
For $f(t)=t^d$, write $r=b/T$, $\zeta=h/T$, and $\vartheta=M/T$ for any $M>0$. Then
\begin{equation}\label{eq:endpointgain}
 \frac{(d+1)\vartheta}{1-r^{d+1}+(d+1)\zeta(1-r^d)}
 \le\frac{MT^d}{\sum_ja_j(b+x_j)^d}
 \le\frac{(d+1)\vartheta}{1-r^{d+1}}.
\end{equation}
\end{lemma}
\begin{proof}
The error on the $j$th interval is nonnegative and at most
$a_j[f(b+x_j)-f(b+x_{j-1})]$. Replace $a_j$ by $h$ and telescope. Integrating $t^d$ and dividing gives~\eqref{eq:endpointgain}.
\end{proof}
Thus $dh/T$ controls the discretization loss. The endpoint coefficient $M$ and the sampled mass $Q$ need not agree. This distinction permits a short, fine window to preserve an amplification approaching $d$.
For a nonnegative-coefficient polynomial $f(t)=\sum_{k=0}^dc_kt^k$ with $f(T)>0$, its exact normalized area is
\begin{equation}\label{eq:mixturearea}
 \frac{\int_b^Tf(t)\,dt}{Tf(T)}
 =\sum_{k=0}^d\lambda_k\frac{1-r^{k+1}}{k+1},\qquad
 \lambda_k=\frac{c_kT^k}{f(T)},\quad\sum_k\lambda_k=1.
\end{equation}
This follows by integrating term by term. A degree bound alone therefore does not force a degree-sized gain: the endpoint may be dominated by lower-degree terms.

\subsection{Equal-weight windows and a local threshold}
Fix an integer $d\ge1$ and normalize a heavy weight to one. Let a window contain an integer $K\ge1$ of light weights $1/R$, and put $m=K/R$ and
\[
 P=(1+m)^d,\quad
 X=\sum_{j=1}^K(1+j/R)^d,\quad
 Y=\sum_{j=1}^K(j/R)^d.
\]
The scalar relevant to propagation of a target factor $\alpha$ is
\begin{equation}\label{eq:gamma}
 \Gamma=\frac{\alpha X-Y}{R(P-\alpha)}\qquad(P>\alpha).
\end{equation}
This expression will be instantiated by actual resource incidences in Appendix~\ref{app:construction}; at present it is a numerical contract. Define
\[
 A_d(m)=\frac{(1+m)^{d+1}-1}{d+1},\qquad
 B_d(m)=\frac{m^{d+1}}{d+1}.
\]
Lemma~\ref{lem:area} gives the finite bounds
\begin{align}
 A_d(m)&\le X/R\le A_d(m)+(P-1)/R,\label{eq:Xarea}\\
 B_d(m)&\le Y/R\le B_d(m)+m^d/R.\label{eq:Yarea}
\end{align}
\begin{proposition}[Finite window certificate]\label{prop:window}
If
\[
 D=P+B_d(m)-\alpha[1+A_d(m)]>0,\qquad
 R>\frac{\alpha(P-1)}D,\qquad mR\in\N,
\]
then $P>\alpha$, $\alpha X-Y>0$, and
\begin{equation}\label{eq:windowmargin}
 1-\Gamma\ge\frac{D-\alpha(P-1)/R}{P-\alpha}>0.
\end{equation}
For fixed $d,m$, the supremum of these finite-grid local thresholds is
\begin{equation}\label{eq:areathreshold}
 \alpha_{\rm area}(d,m)
 =\frac{P+B_d(m)}{1+A_d(m)}
 =\frac{(d+1)(1+m)^d+m^{d+1}}{d+(1+m)^{d+1}}.
\end{equation}
\end{proposition}
\begin{proof}
Since $A_d(m)>B_d(m)$ and $P>1$, the ratio in~\eqref{eq:areathreshold} is less than $P$; hence $D>0$ implies $P>\alpha$. Termwise $X>Y$, so $\alpha X-Y>0$. Subtract the upper estimate obtained from~\eqref{eq:Xarea} and the lower estimate from~\eqref{eq:Yarea} from one to obtain~\eqref{eq:windowmargin}.

For completeness, the finite-grid threshold is $(P+Y/R)/(1+X/R)$. The increasing function $(1+t)^d-t^d$ has a nonnegative right-sum error, so $0\le Y/R-B_d(m)\le X/R-A_d(m)$. Also $P+B_d(m)\ge1+A_d(m)$: their difference is zero at $m=0$, and its derivative is $d(1+m)^{d-1}-[(1+m)^d-m^d]\ge0$. The last inequality follows by integrating the derivative of $t^d$ over $[m,m+1]$. Adding no more error to the numerator than to the denominator cannot increase a ratio at least one. Thus every finite threshold is at most~\eqref{eq:areathreshold}; the two sums converge to their areas as $R\to\infty$ along $mR\in\N$.
\end{proof}
Dropping the nonnegative term $m^{d+1}$ in the numerator produces an explicitly optimized lower envelope. For $d>1$, differentiation gives
\begin{equation}\label{eq:localopt}
 \max_{x>1}\frac{(d+1)x^d}{d+x^{d+1}}
 =d^{(d-1)/(d+1)}>d-2\log d,
 \qquad x=d^{2/(d+1)}.
\end{equation}
Indeed the derivative has the sign of $d^2-x^{d+1}$, and $e^{-u}>1-u$ gives the strict final inequality. Every strictly smaller target has finite rational window parameters, by continuity followed by Proposition~\ref{prop:window}. This is a local propagation threshold, not a statement that all states of a game are unstable. A later propagation may require $\kappa\Gamma<1$ for an additional amplitude factor $\kappa>1$; the two conditions should not be conflated.

\subsection{Ordered marginals and compatible player scales}\label{sec:scales}
Throughout this subsection $d\ge2$. The cases $d=0,1$ are handled separately in Section~\ref{sec:depth}. The integral resource potential $F_e(U)=\int_0^{\sum_{j\in U}w_j}\widehat c_e(z)\,dz$ depends only on load. An ordered right sum also records how individual weights reach that load. For a current user ordering $v_1,\ldots,v_n$, write
\begin{equation}\label{eq:orderedsum}
 P_c=\sum_{j=1}^nv_jc(v_1+\cdots+v_j).
\end{equation}
The order is a rule for defining a set function, not an ordering of actual moves.

\begin{lemma}[Ordered insertion bounds]\label{lem:ordered}
Insert a player of weight $w$ into a fixed ordering, leaving prefix weight $X$ and subsequent weights $v_1,\ldots,v_m$. Let $Y=X+w+\sum_jv_j$. For a nonnegative-coefficient polynomial $c$ of degree at most $d\ge1$,
\[
 0\le\Delta P_c\le dwc(Y).
\]
If all subsequent users have weight at most $w$, the upper bound improves to $wc(Y)$. If all have weight at least $w$, the lower bound improves to $wc(Y)$. A last-position insertion has marginal exactly $wc(Y)$.
\end{lemma}
\begin{proof}
Let $z_j=X+w+\sum_{q\le j}v_q$. The exact marginal is
\[
 \Delta P_c=wc(X+w)+\sum_jv_j[c(z_j)-c(z_j-w)]\ge0.
\]
For $c(z)=z^k$, $1\le k\le d$, each difference is at most $kwY^{k-1}$. On dividing by $wY^k$, the upper bound becomes $u^k+k(1-u)\le k$, where $u=(X+w)/Y$. A constant monomial has exact marginal $w$ times its coefficient, so the bound follows by addition.

For adjacent weights $a\ge b$ and a common endpoint $z$, the contribution with $a$ before $b$ minus that with $b$ before $a$ is
\begin{equation}\label{eq:swap}
 ab\left[\frac{c(z)-c(z-a)}a-\frac{c(z)-c(z-b)}b\right]\le0.
\end{equation}
The inequality holds because the slope of a convex function on a shorter interval with the same right endpoint is larger. Move the inserted player from the end of its subsequent block to its assigned position. If those subsequent weights are all at most $w$, the swaps can only reduce its added contribution; if all are at least $w$, they can only increase it. Appending contributes $wc(Y)$, proving the refinements.
\end{proof}

\subsubsection{Two orders on a resource with two weight types}
We now use the ordering estimate to choose local marginals that can be added across resources.
Suppose the resource has potential optional weights $a>b>0$, with current counts $m,n$. Define the heavy-first and light-first sums
\begin{align}
 D_c(m,n)&=a\sum_{j=1}^mc(ja)+b\sum_{k=1}^nc(ma+kb),\label{eq:D}\\
 U_c(m,n)&=b\sum_{k=1}^nc(kb)+a\sum_{j=1}^mc(nb+ja).\label{eq:U}
\end{align}
Equal-weight users can be freely permuted without changing the value. Applying Lemma~\ref{lem:ordered}, with $Y$ always the post-insertion load, gives
\begin{equation}\label{eq:four}
 0\le\Delta_aD_c\le ac(Y),\quad\Delta_aU_c=ac(Y),\quad
 \Delta_bD_c=bc(Y),\quad bc(Y)\le\Delta_bU_c\le dbc(Y).
\end{equation}
For $r\ge1$, set
\[
 F_{c,r}=(1-1/r)D_c+(1/r)U_c,\qquad
 K_d^{\rm ord}(r)=\max\{r,1+(d-1)/r\}.
\]
Then
\begin{equation}\label{eq:blendmarginals}
 \frac1r ac(Y)\le\Delta_aF_{c,r}\le ac(Y),\qquad
 bc(Y)\le\Delta_bF_{c,r}\le[1+(d-1)/r]bc(Y).
\end{equation}
The minimum of $K_d^{\rm ord}$ is
\begin{equation}\label{eq:sqrtfactor}
 a_d=\frac{1+\sqrt{4d-3}}2,\qquad a_d(a_d-1)=d-1.
\end{equation}
The issue is whether the resource's heavy/light scale ratio $r$ can be made consistent on the other resources of the same players.

\begin{proposition}[Fixed local scales]\label{prop:fixedscales}
Suppose a fixed local resource function has marginal bounds with type coefficients $p_{e,v}>0$, of the form $p_{e,v}v\widehat c_e(Y)\le\Delta_vH_e\le\beta p_{e,v}v\widehat c_e(Y)$. On the multigraph of two-type resources, orient each edge $u\to v$ arbitrarily and label it $q_e=p_{e,v}/p_{e,u}$. Positive resource multipliers $\kappa_e$ and common type scales $t_v$ satisfying $\kappa_ep_{e,v}=t_v$ exist if and only if the oriented product of labels around every closed walk is one.
\end{proposition}
\begin{proof}
The equations imply $t_v/t_u=q_e$, so necessity follows by cancellation around a walk. Conversely choose one root scale in each component and define the others by multiplying labels along a path, using inverse labels backwards. The closed-walk condition makes this independent of the path. Set $\kappa_e=t_u/p_{e,u}$. Single-type resources impose no further condition.
\end{proof}
This is an equivalence for positive multiples of the specified local functions. It does not classify all resource potentials.

\begin{theorem}[Compatible type scales]\label{thm:typescales}
Assume every resource has at most two optional weight types. Form a simple type graph by joining types that share a nonzero-cost two-type resource, oriented from heavy $a$ to light $b$. If positive scales $t_v$ satisfy, on every such edge,
\begin{equation}\label{eq:uniformcontract}
 r_{ab}=t_b/t_a\ge1,\qquad
 r_{ab}\le\beta,\qquad 1+(d-1)/r_{ab}\le\beta,
\end{equation}
then a $\beta$-$\PNE$ exists and every strict $\beta$ improvement terminates.
\end{theorem}
\begin{proof}
Use $t_bF_{\widehat c_e,r_{ab}}$ on a two-type resource. By~\eqref{eq:blendmarginals}, its common lower scales are $at_a$ and $bt_b$, and the corresponding upper scales are at most $\beta$ times these. On a single-type resource of type $v$ use $t_vv\sum_{j=1}^n\widehat c_e(jv)$. The insertion marginal is exactly $t_vv\widehat c_e(Y)$. Apply Lemma~\ref{lem:glue} with $q_i=w_it_{w_i}$.
\end{proof}
\begin{corollary}[Type forests]\label{cor:forest}
If the type graph is a forest, an $a_d$-$\PNE$ exists and every strict $a_d$ improvement terminates, independently of the number of types or their weight ratios.
\end{corollary}
\begin{proof}
Choose a root scale on each tree and impose $t_b=a_dt_a$ along every heavy-to-light edge. Unique paths make these equations consistent, and~\eqref{eq:sqrtfactor} verifies~\eqref{eq:uniformcontract}.
\end{proof}
A long decreasing chain alone therefore does not force a linear approximation in this class. The type graph is distinct from the player--resource incidence graph: two players sharing two resources already give a cycle in the latter and a single edge in the former.

\subsubsection{Cycle constraints and exact contract thresholds}
The forest construction extends to a general type graph precisely when its scale intervals are consistent around cycles.
For $1<\beta<d$, the uniform contract is
\begin{equation}\label{eq:ratiointerval}
 L\le t_b/t_a\le\beta,\qquad L=(d-1)/(\beta-1)>1.
\end{equation}
Writing $z_v=\log t_v$ turns it into difference constraints with arcs $a\to b$ of length $\log\beta$ and $b\to a$ of length $-\log L$.
\begin{theorem}[Cycle imbalance]\label{thm:cycles}
For every undirected simple cycle $C$, let $p_C,q_C$ count heavy-to-light and light-to-heavy steps along a traversal. Both are positive. Define
\[
 \rho_G=\max\left(\{1\}\cup
 \left\{\max(p_C/q_C,q_C/p_C):C\text{ a simple cycle}\right\}\right).
\]
If the type graph has an edge, the least $\beta$ satisfying~\eqref{eq:uniformcontract} is the unique root $\beta_G\in(1,d)$ of
\begin{equation}\label{eq:betarho}
 \beta_G(\beta_G-1)^{\rho_G}=(d-1)^{\rho_G}.
\end{equation}
It is attained, $a_d\le\beta_G<d$, and equality $\beta_G=a_d$ holds exactly when $\rho_G=1$. Equivalently, for $1<\beta<d$ the contract is feasible precisely when $\beta\ge a_d$ and every simple cycle satisfies
\begin{equation}\label{eq:cycleproducts}
 \beta^{q_C}(\beta-1)^{p_C}\ge(d-1)^{p_C},\qquad
 \beta^{p_C}(\beta-1)^{q_C}\ge(d-1)^{q_C}.
\end{equation}
\end{theorem}
\begin{proof}
Summing log-ratio intervals around a cycle requires
$p_C\log L\le q_C\log\beta$ and $q_C\log L\le p_C\log\beta$, giving necessity. For sufficiency, a simple directed cycle of the difference-constraint graph is either a two-arc return on one edge, whose length is $\log\beta-\log L$, or a traversal of an undirected simple cycle. The single-edge condition and~\eqref{eq:cycleproducts} make every such length nonnegative. Hence no negative cycle exists. Add a source with zero arcs; shortest-path distances satisfy all upper difference bounds and supply $z_v$. Exponentiating supplies the scales.

The cycle and single-edge conditions combine as $\log\beta\ge\rho_G\log((d-1)/(\beta-1))$. The difference of these sides is strictly increasing on $(1,d)$, tends to minus infinity at one, and is positive at $d$. Its unique zero is~\eqref{eq:betarho}. At $a_d$ both logarithms agree and are positive, yielding the comparison and its equality criterion. For $\beta\ge d$, constant scales work directly.
\end{proof}
These conditions provide usable decision certificates: feasible scales certify a potential, and a violated product inequality certifies failure of this contract. Failure does not certify that the game has no $\beta$-$\PNE$. For example, put one independent resource of cost $x^d$ on each edge of any type graph, with dedicated endpoint players choosing between that resource and the empty action. The all-empty state is an exact equilibrium, including when the uniform scale contract fails.

The imbalance $\rho_G$ can be obtained without listing every cycle. With $n$ type vertices, it belongs to $\{1\}\cup\{P/Q:1\le Q\le P\le n-1\}$. For any candidate $R\ge1$, the rational difference system
\[
 1\le h_b-h_a\le R\quad(a\to b)
\]
is feasible exactly when $R\ge\rho_G$, by the same cycle proof. A shortest-path computation with rational arc lengths $R,-1$ gives feasibility or a negative cycle. Searching the polynomial-sized candidate set finds $\rho_G$; this complexity statement concerns the explicit type graph, not equilibrium search. At the endpoint, scales can be expressed as $t_v=L^{h_v}$ with a feasible $h$ for $R=\rho_G$. If $\rho_G=P/Q$, the unique endpoint satisfies the integer polynomial $\beta_G^Q(\beta_G-1)^P=(d-1)^P$, so a finite algebraic representation suffices.

A chain of $g$ decreasing edges plus a shortcut from its heaviest to lightest type has $\rho_G=g$ and contract endpoint $\beta(\beta-1)^g=(d-1)^g$. For fixed $g$, this is asymptotic to $d^{g/(g+1)}$. If $g/\log d\to a\in(0,\infty)$, division by $d$ gives a limit $e^{-1/a}$. To see the latter directly, use $\beta-1=(d-1)/\beta^{1/g}$: the bound $\beta<d$ first gives a positive lower bound on $\beta/d$, and then $\log\beta/\log d\to1$. Deleting the shortcut restores a forest and its square-root guarantee. This is a statement about scale conflict, rather than a nonexistence construction.

\subsubsection{Calibration from actual weights and user counts}\label{sec:intervals}
The bounds $0,d$ in~\eqref{eq:four} may be loose. On a nonzero two-type resource let $n_a,n_b$ be the numbers of optional users and define exact finite extrema
\begin{equation}\label{eq:ABextrema}
 \A_e=\min_{0\le m<n_a,\,0\le n\le n_b}
 \frac{\Delta_aD(m,n)}{a\widehat c_e(Y)},\qquad
 \B_e=\max_{0\le m\le n_a,\,0\le n<n_b}
 \frac{\Delta_bU(m,n)}{b\widehat c_e(Y)}.
\end{equation}
They use all local count sets within the potential population, whether or not every such set occurs in a global state. The dictionary studied here is the nonnegative span $xD+yU$; the certificate requires insertion bounds on that complete local domain.
\begin{theorem}[Finite interval calibration]\label{thm:interval}
One has $0<\A_e\le1\le\B_e\le d$. Fix $\beta\ge1$, normalize the heavy type scale to one, and write $r=t_b/t_a$. There are $x,y\ge0$ giving common marginal bounds at factor $\beta$ if and only if
\begin{equation}\label{eq:fourLP}
 \A_ex+y\ge1,\qquad x+y\le\beta,\qquad
 r\le x+y,\qquad x+\B_ey\le\beta r.
\end{equation}
If nonempty, the projection onto $r$ is a compact interval $[l_e,\beta]$, $l_e>0$, with rational endpoints for rational data and rational $\beta$.

Writing $A=\A_e,B=\B_e$, define
\[
 \lambda_-=
 \begin{cases}0&A=1,\\ \max\{0,(\beta^{-1}-A)/(1-A)\}&A<1,\end{cases}
 \quad
 \lambda_+=
 \begin{cases}1&B=1,\\ \min\{1,(\beta-1)/(B-1)\}&B>1.\end{cases}
\]
The interval is nonempty exactly when $\lambda_-\le\lambda_+$. In that case
\begin{equation}\label{eq:lowinterval}
 l_e=\min_{\lambda\in[\lambda_-,\lambda_+]}
 \frac{1+(B-1)\lambda}{\beta[A+(1-A)\lambda]}.
\end{equation}
The minimum is at $\lambda_-$ when $AB\ge1$ and at $\lambda_+$ when $AB\le1$.
\end{theorem}
\begin{proof}
The heavy insertion in $D$ has a strictly positive own term, and the finite population gives a positive minimum. Its maximum normalized marginal is one, attained with no light users. For light insertion in $U$, the minimum is one, attained with no heavy users. The other two normalized marginals are identically one. Thus the exact extrema of the nonnegative mixture are $Ax+y,x+y$ for heavy users and $x+y,x+By$ for light users; these are exactly~\eqref{eq:fourLP}.

Put $s=x+y$, $\lambda=y/s$. Then $1\le s\le\beta$, and feasibility is
\[
 s[A+(1-A)\lambda]\ge1,\quad
 s[1+(B-1)\lambda]\le\beta r,\quad r\le s\le\beta.
\]
For fixed $\lambda$, existence is equivalent to $A+(1-A)\lambda\ge\beta^{-1}$ and $1+(B-1)\lambda\le\beta$, namely the stated $\lambda$ interval. The least $r$ at that $\lambda$ is the ratio in~\eqref{eq:lowinterval}, obtained by taking $s=[A+(1-A)\lambda]^{-1}$. Its derivative has the sign of $AB-1$. The feasible polytope is convex and bounded in $x,y,r$, with $r\ge1/\beta>0$, so the projection is a compact interval. From any feasible $(x,y,r)$, multiply $x,y$ by $\beta/s$ and set $r=\beta$; the inequalities remain valid. Hence its upper endpoint is $\beta$. All formulas are rational when the input is rational.
\end{proof}
Unlike~\eqref{eq:uniformcontract}, an actual interval can include $r<1$. After combining parallel resources by $l_{ab}=\max_e l_e$, common type scales exist exactly when
\begin{equation}\label{eq:heterogeneous}
 l_{ab}\le t_b/t_a\le\beta
\end{equation}
holds on every edge. The same difference-constraint proof applies, now with reverse arc length $-\log l_{ab}$. Equivalently, every traversal of a simple cycle, with forward-edge set $F$ and reverse-edge set $R$, obeys
\[
 \prod_{e\in F}l_e\le\beta^{|R|},\qquad
 \prod_{e\in R}l_e\le\beta^{|F|}.
\]
These heterogeneous constraints should not be replaced by the uniform imbalance formula~\eqref{eq:betarho}. On a forest, nonempty local intervals are sufficient, by assigning scales along unique paths.

For a resource with $0<A<1<B$, its least feasible dictionary factor is the unique positive root $q_e>1$ of
\begin{equation}\label{eq:quadratic}
 (1-A)q_e^2+(AB-1)q_e-(B-1)=0.
\end{equation}
If $A=1$ or $B=1$, it is one. To prove~\eqref{eq:quadratic}, on $1\le\beta\le\min\{1/A,B\}$ cross-multiply $\lambda_-\le\lambda_+$; the resulting inequality is that the polynomial in~\eqref{eq:quadratic} is nonnegative. It has one positive root, is negative at one, and is nonnegative at both $1/A$ and $B$. Above their minimum feasibility persists by relaxing the same bounds. Consequently the exact minimum dictionary factor on a forest is $\max_eq_e$, with the endpoint attained. Exactness here refers throughout to this dictionary, these type scales, and all local insertions.

Rational calibration can avoid numerical logarithms. Treat difference-constraint arc lengths as logarithms of rational multipliers $\beta$ and $1/l_e$. A layered shortest-path calculation multiplies these rationals and compares the products. At most $n$ layers suffice for feasibility or a cycle certificate; path products then have polynomial bit length in explicit rational input. The extrema~\eqref{eq:ABextrema} use at most $n_a(n_b+1)+(n_a+1)n_b$ count evaluations. With explicitly listed players and strategies and dense polynomial input these certificate calculations have polynomial size. They do not bound the length of a subsequent improvement path.

For example, one optional user of weight three and one of weight one on cost $x^2$ give $A=7/8,B=11/8$. At $\beta=10/9$, the exact interval is $[43/40,10/9]$, with lower-endpoint mixture $x=8/9,y=2/9$. Indeed $Ax+y=1$, $x+y=10/9$, and $x+By=43/36=\beta(43/40)$. This improves the uniform square-root certificate on this resource and any forest of such calibrated resources, without asserting that the game's true best equilibrium factor is $10/9$.

\subsection{Weight depth and the onset of linear approximation factors}\label{sec:depth}
Type-cycle compatibility and weight depth measure different obstructions. The next theorem allows arbitrarily many types per resource and uses a common rank on individual players.

Absorb globally fixed players as background. On the remaining players put an arc $i\to j$ if $w_i>w_j$ and some resource belongs to at least one strategy of each. Let $h$ be the maximum number of vertices on a directed path, and $k$ the number of distinct remaining weights. The graph is acyclic and $h\le k$. In particular, the graph uses potential co-occupancy, not co-occupancy at one chosen state. Aggregating players of the same weight into types can create paths that cannot be followed through the same intermediate player.

For $d\ge2,h\ge1$, let $A_{d,h}$ be the unique positive root
\begin{equation}\label{eq:heightroot}
 A=\sum_{j=0}^{h-1}\left(\frac{d-1}{A}\right)^j.
\end{equation}
It lies in $[1,d)$: the right side minus $A$ is strictly decreasing, is nonnegative at one, and is negative at $d$.
\begin{theorem}[Player-height potential]\label{thm:height}
Every game of degree at most $d\ge2$ and potential decreasing-weight height $h\ge1$ has an $A_{d,h}$-$\PNE$, and every strict improvement by this factor terminates. The same conclusion holds with $k$ in place of $h$.
\end{theorem}
\begin{proof}
Give player $i$ rank $r_i$, the length in vertices of a longest directed path ending there. Potential co-users of unequal weight have strictly increasing ranks as their weights decrease; equal-rank co-users therefore have equal weights. For $b=1,\ldots,h$, order resource users by the rank order
\[
 \pi_b=(h,h-1,\ldots,b+1,1,2,\ldots,b)
\]
and use the right sum~\eqref{eq:orderedsum}, denoted $P_{b,e}$. In a given rank, identical weights can be reordered so that an insertion is last in its rank. For a player of rank $i$, Lemma~\ref{lem:ordered} gives normalized marginal intervals $[1,d]$ when $b<i$, the exact value one when $b=i$, and $[0,1]$ when $b>i$.

Write $A=A_{d,h}$, $t=(d-1)/A$, $\lambda_b=t^{b-1}/A$, and $S_i=\sum_{b\le i}\lambda_b$. Equation~\eqref{eq:heightroot} gives $\sum_b\lambda_b=1$. The mixture $F_e=\sum_b\lambda_bP_{b,e}$ has insertion bounds
\[
 S_iw\widehat c_e(Y)\le\Delta_iF_e
 \le[1+(d-1)S_{i-1}]w\widehat c_e(Y)
 =AS_iw\widehat c_e(Y).
\]
The last identity follows directly from the finite geometric sums, including $t=1$. Thus $q_i=w_iS_{r_i}$ is independent of the resource. Apply Lemma~\ref{lem:glue}. The root is nondecreasing in $h$, so $h\le k$ yields the second assertion.
\end{proof}
For $d=0$, each player independently minimizes its sum of constant latencies. For $d=1$, an exact resource potential is $\ell_e(x_e^2+\sum_{i:e\in s_i}w_i^2)/2+\gamma_ex_e$ when $c_e(x)=\ell_ex+\gamma_e$; its insertion marginal is $w_ic_e(Y)$. Thus these degrees have exact equilibria without a height restriction.

\begin{theorem}[Necessary depth, with its exact inversion]\label{thm:depthnecessary}
For $1\le\alpha<d$, define
\begin{equation}\label{eq:depthinverse}
 H_d(\alpha)=
 \begin{cases}
 \displaystyle\frac{\log(d-\alpha)}{\log((d-1)/\alpha)},&\alpha\ne d-1,\\[5pt]
 d-1,&\alpha=d-1.
 \end{cases}
\end{equation}
A game without an $\alpha$-$\PNE$ must have $h>H_d(\alpha)$ and $k>H_d(\alpha)$. No game without an $\alpha$-$\PNE$ exists for $\alpha\ge d$.
\end{theorem}
\begin{proof}
Theorem~\ref{thm:height} implies $\alpha<A_{d,h}<d$. Let $f_h(x)=\sum_{j=0}^{h-1}((d-1)/x)^j-x$, a strictly decreasing function. With $t=(d-1)/\alpha\ne1$,
\[
 (t-1)f_h(\alpha)=t^h-(d-\alpha).
\]
If $t>1$, $f_h(\alpha)>0$ means $t^h>d-\alpha$; if $t<1$, it means $t^h<d-\alpha$. Taking logarithms in either case gives $h>H_d(\alpha)$. At $t=1$ the condition is $h>d-1$. The inequality for $k$ follows from $h\le k$.
\end{proof}
For each fixed $c\in(0,1)$ this yields
\begin{equation}\label{eq:fixedclowerdepth}
 h,k>\frac{\log((1-c)d)}{\log((1-1/d)/c)}
 =\frac{\log d+\log(1-c)}{\log(1/c)}+o_c(1)
\end{equation}
whenever no $cd$-$\PNE$ exists and $d$ is sufficiently large. It also calibrates the potential's transition: if $h/\log d\to a\in(0,\infty)$, then $A_{d,h}/d\to e^{-1/a}$. Indeed for each fixed $x\in(0,1)$ the ratio $H_d(xd)/\log d$ tends to $1/\log(1/x)$, and $A_{d,h}>xd$ is equivalent to $h>H_d(xd)$. Taking $x$ on either side of $e^{-1/a}$ proves the limit. The endpoint cases give $A_{d,h}=o(d)$ if $h=o(\log d)$ and $A_{d,h}/d\to1$ if $h/\log d\to\infty$.

The fixed-$c$ qualification cannot be dropped. At $\alpha=d-1$, nonexistence already requires $h,k\ge d$. More generally the constants in~\eqref{eq:fixedclowerdepth} diverge as $c$ tends to one. High depth is necessary, not sufficient, for instability: an arbitrarily deep potential co-occupancy graph can coexist with an exact equilibrium.

For $1\le\alpha<d$, there is also a witness interpretation for a cycle of strict $\alpha$ improvements. Retain only players moving in the cycle, absorb the others as background, and draw decreasing-weight arcs only between players actually sharing a resource at some endpoint of that cycle. Apply the proof of Theorem~\ref{thm:height} with ranks from this graph. Its insertion bounds are needed only for the old deletion and new addition endpoints that occur in the cycle, where the ranks are compatible with the weights. If its height were at most $H_d(\alpha)$, the same potential would strictly decrease around the cycle, a contradiction. Thus even such a cycle must collectively witness a longer decreasing-weight path, possibly using different resources and endpoint states for its different edges. This statement does not assert nonexistence from the cycle.

\subsection{Realizing stability on a structured set of states}\label{sec:repair}
The preceding potentials certify deviations throughout a given game. We now consider a complementary construction: first analyze deviations within a structured set of states, then add legal resources that exclude all other states. The added costs are part of the resulting game and must be included in its stability comparisons.

Let a binary base game $G_B$ have ordered classes $\mathcal C_a=((a,0),\ldots,(a,n_a-1))$, with common weight $w_a>0$ inside each class. Let $F_n$ be the binary strings with at most one change between consecutive positions; we call these strings \emph{fronts} and put $\F=\prod_aF_{n_a}$. A \emph{core deviation} is a real single-player flip whose two endpoints both belong to $\F$. This restricted deviation system need not be a normal-form subgame because $\F$ need not be a Cartesian product over individual players.

\subsubsection{Legal gates and the complete stable set}
We realize the front condition using resources shared by three players of the same weight.
For every class and $1\le j<k<n_a$, add two fresh resources used only by players $(a,0),(a,j),(a,k)$. Their resource sides are their action bits xor $(0,1,0)$ respectively. The two sides have the same nonnegative nondecreasing latency $c_a$, which may be a nonnegative-coefficient polynomial. Resource pairs are disjoint from all other pairs and the base game. No players, weights, or action labels are added. Write
\begin{equation}\label{eq:gateparameters}
 L_a=\binom{n_a-1}{2},\quad z_a=c_a(w_a),\quad
 u_a=c_a(2w_a),\quad v_a=c_a(3w_a),
\end{equation}
with $L_a=0$ if $n_a\le2$. A gate is triggered precisely for action patterns $010$ or $101$, when all three players use the same side and each pays $v_a$. Otherwise its users split into a pair and a singleton. The number of gates involving player $(a,i)$ is
\begin{equation}\label{eq:gatecounts}
 r_{a,i}=\begin{cases}
 L_a,&i=0,\\ n_a-2,&i>0,\ n_a\ge3,\\0,&n_a\le2.
 \end{cases}
\end{equation}
In all cases $r_{a,i}\le L_a$. Let $B_a>0$ bound every base-game cost of every player of class $a$ over all states. Such a bound can be obtained by evaluating each resource at the total weight of all its potential users and summing over the resources a player may use.

\begin{theorem}[Stable-core realization]\label{thm:core}
Fix $\alpha\ge1$. Suppose, for each $n_a\ge3$,
\begin{equation}\label{eq:RD}
 v_a>\alpha(B_a+L_au_a).
\end{equation}
In the augmented game $G_+$, the entire set of $\alpha$-$\PNE$ is
\begin{equation}\label{eq:toleranceset}
 \mathcal T_\alpha=\left\{s\in\F:
 C_i^B(s)-\alpha C_i^B(s^{\xor i})
 \le(\alpha-1)r_{a,i}u_a
 \quad\text{for every core flip by }i\in\mathcal C_a\right\}.
\end{equation}
At $\alpha=1$ this is exactly the base core's stable set. For $\alpha>1$ the displayed tolerance is essential.
\end{theorem}
\begin{proof}
All gates of a class are untriggered exactly for front strings. A front has no alternating subsequence of the form $b,1-b,b$. If a string is not a front, let $i>0$ be its first bit different from its first bit $b$; some later $j$ equals $b$, so gate $(0,i,j)$ triggers.

At a state outside $\F$, perform this single-player correction $i\mapsto b$. Its old gate cost is at least $v_a$. Every gate involving it is now untriggered: when $i$ is the middle index it equals the anchor; when it is the last index, the middle index is earlier and already equals the anchor. The new base cost is at most $B_a$ and new gate cost at most $L_au_a$. Thus the correction has strictly positive $\alpha$ margin by~\eqref{eq:RD}. No outside state is stable.

Starting inside $\F$, a flip leaving $\F$ triggers a gate involving the mover. Its new cost is at least $v_a$, whereas its old cost is at most $B_a+L_au_a<v_a/\alpha$. Such a flip cannot be a strict $\alpha$ improvement.

Finally take a core flip. In each gate involving the mover, both endpoint patterns are untriggered. The mover must be in a two-player group before and after: moving from a singleton would join the other two and trigger the gate. Its cost is therefore exactly $u_a$ on both endpoints of every participating gate. Its full margin is
\begin{equation}\label{eq:coremargin}
 C_i^+(s)-\alpha C_i^+(s^{\xor i})
 =C_i^B(s)-\alpha C_i^B(s^{\xor i})-(\alpha-1)r_{a,i}u_a.
\end{equation}
The three cases exhaust every state and legal deviation, proving the formula.
\end{proof}
Every $\alpha$-stable base core state lifts, but approximate stable states can also be added: the common gate baseline absorbs a small positive base margin. Exact stability is preserved in both directions. In particular, a core with one exact stable state cannot be made completely unstable merely by increasing these gate coefficients. Conversely, all-state nonexistence follows if every core state has a preserving exit exceeding its corresponding tolerance.

The nonzero baseline reflects an obstruction for added resources with nonnegative-coefficient polynomial latencies. If every action of a player appears in some core state, an added cost of this kind that vanishes for that player on all core states must vanish on all its states. Indeed a nonzero nonnegative-coefficient resource has positive cost at any positive load, including the player's own weight; hence none of its actions can contain such a resource. One cannot assume zero gate leakage on a full-projection core without proving a different implementation.

\subsubsection{Repair paths and the degree budget}
The same construction gives an explicit correction path and a condition for implementing the gates at a prescribed degree.
The correction in the proof lowers by one the number of nonanchor bits differing from the anchor. Initially that number is at most $n_a-2$ in an illegal class. The specified correction schedule therefore reaches $\F$ in at most
\begin{equation}\label{eq:repairsteps}
 \sum_a\max\{n_a-2,0\}
\end{equation}
steps, each a strict $\alpha$ improvement under~\eqref{eq:RD}. This does not imply termination of arbitrary improvement paths after entering the core.

For a monomial implementation, let $B=\max_aB_a$, $L=\max_aL_a\ge1$, and choose
\[
 M=2\alpha B,\qquad c_a(x)=M(x/(3w_a))^D,\qquad\zeta_D=(2/3)^D.
\]
Then $v_a=M$ and $u_a=M\zeta_D$. The condition $\zeta_D\le1/(4\alpha L)$ implies~\eqref{eq:RD}, with margin at least $\alpha B/2$. If every core state has a preserving base exit of margin at least $\eta>0$, it is enough in addition that
\begin{equation}\label{eq:compilerbudget}
 (\alpha-1)LM\zeta_D\le\eta/2.
\end{equation}
The resulting game has no $\alpha$-$\PNE$. A convenient stronger sufficient budget, when $\eta\le B$, is
\begin{equation}\label{eq:safebudget}
 (2/3)^D\le\frac{\eta}{4\alpha^2BL}.
\end{equation}
The output degree is $\max\{p,D\}$ when the base degree is at most $p$. The number of added resources is $2\sum_aL_a$, and the number of new player--action--resource incidences is $6\sum_aL_a$. The coefficient $M/(3w_a)^D$ is rational for rational data and has bit length $O(\operatorname{bits}(M)+D(1+\operatorname{bits}(w_a)))$. Thus keeping a prescribed degree requires a real inequality proving $D$ is within budget; it does not follow from finite realizability.

This design can also be expressed as a finite linear certificate. For a base cost function $B_i(s)$ and already legal disjoint modules with player costs $H_i^j(s)$, nonnegative coefficients $\lambda_j$ yield $C_i^\lambda(s)=B_i(s)+\sum_j\lambda_jH_i^j(s)$. For a deviation $s\to t$ by player $i$, write $\Delta_{\alpha,i}X(s,t)=X_i(s)-\alpha X_i(t)$ for any player-cost family $X$. Chosen exits from states to exclude impose strict linear inequalities
\[
 \Delta_{\alpha,i}B(s,t)+\sum_j\lambda_j\Delta_{\alpha,i}H^j(s,t)>0,
\]
and all deviations of states to retain impose the corresponding weak reverse inequalities. A common slack variable $0\le\rho\le1$ replacing strict positivity by $\ge\rho$ certifies feasibility with $\rho>0$. Rational feasible data admit rational realizations. This is a finite template-selection method, possibly involving all states; it cannot convert arbitrary prescribed cost tables into legal congestion resources, nor can an infeasible template rule out all other realizations.

\section{All-state counterexamples from finite windows}\label{sec:application}
The endpoint estimates and the realization theorem now give a family with a strict improvement at every state. Adjacent windows transmit the improvement between weight classes; a return window covers the phase configurations left by the adjacent tests. The anchored gates from Theorem~\ref{thm:core} then cover all states outside the product of fronts. We give the actual parameters and resource sets before proving the asymptotic consequence.

\subsection{A family with explicit rational parameters}
The construction has one parameter controlling the target factor and another controlling the polynomial degree. All resource incidences and coefficients are specified by the following finite formulas.
\begin{theorem}[Uniform family]\label{thm:uniformfamily}
Let integers $n\ge2$ and $d\ge16$ satisfy
\begin{equation}\label{eq:uniformconditions}
 d\ge8n,\qquad T=(1+1/(8n))^d\ge16nd,\qquad
 d\log(3/2)>200n(\log d)^2.
\end{equation}
There is a finite rational game with positive pure degree-$d$ monomial latencies, two nonempty disjoint strategies per player, and a strict $(1-1/n)d$ improvement at every state.
\end{theorem}
We specify the game and then deduce the theorem from the estimates in Appendix~\ref{app:construction}. Define
\begin{equation}\label{eq:appfirstparams}
 \delta=1/n,\quad\alpha=(1-\delta)d,\quad
 \theta=\delta/8,\quad\kappa=1+\theta,\quad
 b=1-\delta/2,\quad K=d^4,\quad R=8nd^4.
\end{equation}
Thus $K/R=\theta$ and the atomic resolution is $1/R$; these two quantities have separate roles. Set
\begin{align}
 H&=T-\alpha,&
 X&=\sum_{r=1}^K(1+r/R)^d,&
 Y&=\sum_{r=1}^K(r/R)^d,\label{eq:appXY}\\
 J&=\alpha X-Y,&S&=\kappa J/H.\nonumber
\end{align}
Let $g\ge2$ be the least integer with $b^g\le(192d^2)^{-1}$, and put
\begin{equation}\label{eq:appdepthparams}
 N=R^g/d,\qquad\varepsilon=R^{-g},\qquad N\varepsilon=1/d.
\end{equation}
The number $N$ is an integer because $d$ divides $R$. Define the return sums and amplitudes
\begin{align}
 W_0&=(1+1/d)^d,&
 U&=\sum_{r=1}^N(1+r\varepsilon)^d,&
 V&=\sum_{r=1}^N(r\varepsilon)^d,&Z&=U-\alpha V,\label{eq:appreturn}\\
 A_a&=S^a\quad(0\le a<g),&&
 A_R&=H/(4\alpha W_0).\nonumber
\end{align}
All these definitions use only integer arithmetic, finite rational sums and powers, and rational comparisons. Logarithms enter the estimates but not the game data.

\subsubsection{Players, resource pairs, and actions}\label{app:resources}
The parameters determine a game directly: each player selects one side of every resource pair in which it participates.
There are players $P_{a,j}$ with $0\le a\le g$ and $j\in\mathbb Z/N\mathbb Z$, of weight $w_a=R^{-a}$. Use representatives $j=0,\ldots,N-1$ and the global position $\pos(a,j)=j(g+1)+a$.

A test $(a,b,k,o)$, $a<b$, has a pair of fresh resources $e_{\tau,j,0},e_{\tau,j,1}$ for each anchor $h=P_{a,j}$. The anchor's action bit selects its resource side. Its window users are $u=P_{b,j+t\bmod N}$ for $t=0,\ldots,k-1$. A window user with action bit $z_u$ selects side
\begin{equation}\label{eq:appincidence}
 z_u\xor\ind[\pos(u)>\pos(h)]\xor\ind[o=-1].
\end{equation}
No other user occupies the pair. Use the $g$ adjacent tests $(a,a+1,K,-1)$ with both resource latencies $(A_a/w_a^d)x^d$, and the return test $(0,g,N,+1)$ with both latencies $A_Rx^d$. All pairs in all tests are distinct. A player's action is the union of the sides prescribed by its bit over all pairs in which it participates. Every player participates, so both actions are nonempty; the actions are disjoint since each pair contributes opposite sides. There are $(g+1)N$ players and $2(g+1)N$ base resources, all with positive rational data.


For the full game, set
\begin{equation}\label{eq:maincompletion}
 B_*=(g+1)R^gT,\qquad M=2\alpha B_*.
\end{equation}
For every class $a$ and every $1\le j<k<N$, add the anchored pair of Section~\ref{sec:repair} on $P_{a,0},P_{a,j},P_{a,k}$, with latency $M(x/(3w_a))^d$ on both sides. These are ordinary congestion resources. They add no players and preserve the two disjoint actions of every player.

\begin{proof}[Proof of Theorem~\ref{thm:uniformfamily}]
Lemma~\ref{lem:appestimates} gives the endpoint, propagation, and return estimates for these parameters. The phase analysis in Appendix~\ref{app:frontcoverage} shows that every product-front state has a preserving single-player deviation with base margin at least $\eta=\theta H/4$. Appendix~\ref{app:samedegree} proves that $B_*$ bounds all base costs and verifies the degree-$d$ gate budget~\eqref{eq:safebudget}. Theorem~\ref{thm:core} therefore excludes every nonfront state, while each front exit retains margin at least $\eta/2$. Thus every state has a strict $\alpha=(1-1/n)d$ improvement. The displayed resource formulas give positive rational data and two nonempty disjoint actions for every player.
\end{proof}

For a single fully specified member of the family, take $n=2$ and $d=2^{21}$ in these formulas. Corollary~\ref{cor:fixedn} verifies the parameter conditions, so the resulting rational game has no $2^{20}$-$\PNE$. This choice provides an explicit instance without optimizing its number of players or resources. Varying $n$ and $d$ gives the stronger asymptotic statement below.

\subsection{The near-degree bound and the universal threshold}
We now choose the target factor as a function of the degree. The uniform conditions of Theorem~\ref{thm:uniformfamily} allow this choice to approach $d$ while retaining the same legal resource construction.

\begin{theorem}[All-state near-degree lower bound]\label{thm:near}
For every integer $d\ge2^{20}$ there is a finite weighted congestion game with positive rational weights and latencies $c_e(x)=a_ex^d$, $a_e>0$ rational, such that every player has exactly two nonempty disjoint resource-set strategies and
\begin{equation}\label{eq:nearlower}
 \forall s\ \exists i,t_i:\quad
 C_i(s)>\bigl(d-1200\lceil\log_2d\rceil^2\bigr)C_i(t_i,s_{-i}).
\end{equation}
\end{theorem}
\begin{proof}
Put $\ell=\lceil\log_2d\rceil$ and $n=\lceil d/(1200\ell^2)\rceil$. Appendix~\ref{app:tail} checks all three conditions in~\eqref{eq:uniformconditions} for every $d\ge2^{20}$. Theorem~\ref{thm:uniformfamily} gives a strict improvement at factor $d-d/n\ge d-1200\ell^2$. Nonnegative costs give the claimed margin at the smaller factor.
\end{proof}
The parameter check is uniform in $n$ and $d$. This is needed because $n$ now grows with the degree; a statement with a separate unspecified threshold for each fixed $n$ would not imply Theorem~\ref{thm:near}.

Let $a_*(d)$ be the infimum of $\alpha\ge1$ such that every finite game of degree at most $d$ has an $\alpha$-$\PNE$. Caragiannis and Fanelli's monomial potential~\cite[Theorem~2]{CF2021}, extended termwise as in the Introduction, and Theorem~\ref{thm:near} give
\begin{equation}\label{eq:universalthreshold}
 d-1200\lceil\log_2d\rceil^2\le a_*(d)\le d,
 \qquad \lim_{d\to\infty}\frac{a_*(d)}d=1.
\end{equation}
The infimum formulation needs no claim about an extremal finite game at the precise universal threshold.

There is also a matching structural consequence. For fixed $c\in(0,1)$ define $H_c(d)$ and $K_c(d)$ as the least possible potential player height and the least possible number of distinct movable weights among degree-at-most-$d$ games with no $cd$-$\PNE$, with value $+\infty$ if none exist.
\begin{corollary}[Minimum depth for a fixed linear obstruction]\label{cor:mindepth}
For every fixed $c\in(0,1)$, both $H_c(d)$ and $K_c(d)$ are $\Theta_c(\log d)$ for all sufficiently large integer $d$.
\end{corollary}
\begin{proof}
The lower bound is~\eqref{eq:fixedclowerdepth}. For the upper bound choose an integer $n>1/(1-c)$ once and use the uniform construction in Section~\ref{sec:application}, verified in Appendix~\ref{app:construction} with target $(1-1/n)d$. Its sufficient conditions hold for every large enough $d$, as shown in Corollary~\ref{cor:fixedn}. The players have weights $R^{-a}$ for $a=0,\ldots,g$, with $g\le13n\log d$. Repair adds no players or weights. Hence $k=g+1$. Moreover the particular players $P_{0,0},\ldots,P_{g,0}$ form a potential decreasing-weight path: the adjacent test anchored at clone zero includes light clone zero in its window. Thus $h=g+1$ as well. The constants may depend on the fixed $c$ through $n$.
\end{proof}
The lower and upper bounds match in order, not in leading coefficient. In the varying near-one regime of Theorem~\ref{thm:near}, the chosen depth is larger and this fixed-$c$ statement does not assert a uniform logarithmic bound.

\section{Applications to stable states and social optimization}\label{sec:positive}
The counterexamples use large endpoint amplification and enough weight levels to propagate it. We now apply the area comparison to small actual shares and use the scale and gate constructions to study stable states with social-cost guarantees.

\subsection{Existence under small atomic shares}
To quantify the effect of small actual atomic shares, define, for $0\le u\le1$,
\begin{equation}\label{eq:kfunction}
 K_d(u)=\int_0^1(1-uv)^d\,dv
 =\frac{1-(1-u)^{d+1}}{(d+1)u},\qquad K_d(0)=1.
\end{equation}
It is nonincreasing in $u$, and $K_k(u)\ge K_d(u)$ when $k\le d$.
\begin{theorem}[Small-share existence]\label{thm:smallatom}
Suppose that, at every feasible state where an optional player $i$ occupies resource $e$, its share of the actual total load obeys $w_i/x_e(s)\le\eta\le1$. Then a $K_d(\eta)^{-1}$-$\PNE$ exists, and every strict improvement by that factor terminates.
\end{theorem}
\begin{proof}
Use $\Phi(s)=\sum_e\int_{b_e}^{x_e(s)}c_e(t)\,dt$. The contribution removed by a player of weight $w$ from total load $L$ is, for each monomial, exactly $wL^kK_k(w/L)$. Thus it is at least $wK_d(\eta)c_e(L)$. On an added resource monotonicity bounds the contribution by $wc_e(L')$. Retained costs do not change, so
\[
 \Phi(t)-\Phi(s)\le w_i[C_i(t)-K_d(\eta)C_i(s)].
\]
In this sum the retained resources add the nonnegative term $w_i(1-K_d(\eta))C_{i,H}$ to the right side. A strict $K_d(\eta)^{-1}$ improvement decreases this finite-state potential.
\end{proof}
For fixed $a>0$, $K_d(a/d)^{-1}\to a/(1-e^{-a})$; when $d\eta\to0$ the guarantee tends to one. The hypothesis refers to occupied loads at all relevant states. For the resourcewise preprocessing of Section~\ref{sec:model}, the largest optional share on resource $e$ is exactly $w_i/(b_e+w_i)$: every other optional user can independently choose a legal strategy omitting $e$, while $i$ chooses one containing it. Thus explicit strategy menus make this bound directly checkable. In particular a resource with no compulsory load offers only the bound one under this all-state condition. Adding many potential users, or checking only one state, does not establish a small-share guarantee. This positive theorem and the endpoint gain describe the same analytic phenomenon from opposite atomic scales.

\subsection{Choosing a stable state with a social-cost certificate}\label{sec:regularization}
The same potentials also give a tradeoff between approximate stability and social cost. Here assume no compulsory users on any resource, so that the load in the potential is the entire load in~\eqref{eq:welfare}. Let $\Phi$ be a potential from Theorem~\ref{thm:typescales} with factor $\beta\le d$, resource multiplier $\kappa_e=t_{\rm light}$ on two-type resources and $\kappa_e=t_v$ on single-type resources. Ignore resources without potential users, and assume at least one remains. Put
\[
 m_0=\frac{\min_e\kappa_e}{d+1},\qquad
 M_0=\max_e\kappa_e,\qquad t_{\min}=\min_vt_v.
\]
Every ordered right sum at load $X$ lies between $\int_0^Xc_e(z)\,dz$ and $Xc_e(X)$; the integral is at least $Xc_e(X)/(d+1)$, term by term. Convex combinations preserve this, so
\begin{equation}\label{eq:costsandwich}
 m_0W(s)\le\Phi(s)\le M_0W(s).
\end{equation}
\begin{proposition}[Stability and cost regularization]\label{prop:regularization}
For $\eta\ge0$, any global minimizer $s_\eta$ of $\Phi+\eta W$ satisfies
\begin{equation}\label{eq:regguarantees}
 s_\eta\text{ is a }\gamma_\eta\text{-}\PNE,\qquad
 \gamma_\eta=\beta+\frac{(d+1-\beta)\eta}{t_{\min}+\eta},\qquad
 W(s_\eta)\le\frac{M_0+\eta}{m_0+\eta}\min_sW(s).
\end{equation}
\end{proposition}
\begin{proof}
For an inserted weight $w$ ending at load $Y$, the marginal of $Xc_e(X)$ lies between $wc_e(Y)$ and $(d+1)wc_e(Y)$. For a monomial this is the identity $Y^{k+1}-(Y-w)^{k+1}=w\sum_{r=0}^kY^{k-r}(Y-w)^r$, whose sum lies between $Y^k$ and $(k+1)Y^k$. Hence the new lower and upper marginal coefficients for player $i$ are $w_i(t_{w_i}+\eta)$ and $w_i(\beta t_{w_i}+(d+1)\eta)$. Their ratio is at most $\gamma_\eta$, maximized at $t_{\min}$. Lemma~\ref{lem:glue} gives stability. Global minimality and~\eqref{eq:costsandwich} give the cost comparison.
\end{proof}
The stability conclusion uses only local minimality, whereas the social-cost comparison uses global minimality. With compulsory background loads, the target $W$ also includes their varying costs, so~\eqref{eq:costsandwich} requires a separate baseline estimate. The bound is informative when the resulting scale range is controlled.

\subsection{Social compression and transfer of quality}\label{sec:social}
The gate construction gives more than stable-set localization. Here $C_i^B,W^B$ denote player costs and social cost in the base game $G_B$, and $C_i^+,W^+$ denote those in the augmented game $G_+$. Let $\nu_a(s)$ count triggered gates of class $a$, and put
\begin{equation}\label{eq:socialgate}
 d_a=3v_a-2u_a-z_a,\qquad
 K_0=\sum_aw_aL_a(2u_a+z_a).
\end{equation}
The symbol $d_a$ is a social penalty, distinct from the polynomial degree $d$.
\begin{lemma}[Exact social decomposition]\label{lem:socialdecomposition}
At every state,
\begin{equation}\label{eq:allsocial}
 W^+(s)=W^B(s)+K_0+\sum_aw_a\nu_a(s)d_a.
\end{equation}
In particular $W^+(s)=W^B(s)+K_0$ throughout $\F$.
\end{lemma}
\begin{proof}
An untriggered gate has two users paying $u_a$ and one paying $z_a$, while a triggered gate has three users paying $v_a$. The users have the same weight $w_a$. Summing the two possible contributions gives the identity.
\end{proof}
The same front constant-shift identity holds for any objective $J_\gamma=\sum_a\gamma_a\sum_{i\in\mathcal C_a}C_i$ with coefficients constant within a class, replacing $w_a$ by $\gamma_a$. It need not hold for arbitrary player coefficients or the maximum player cost.

Under~\eqref{eq:RD}, if the stable set is nonempty, Theorem~\ref{thm:core} and Lemma~\ref{lem:socialdecomposition} imply the exact stable selection identity
\begin{equation}\label{eq:stableselection}
 \min_{s\in\PNE_\alpha(G_+)}W^+(s)
 =K_0+\min_{s\in\mathcal T_\alpha}W^B(s).
\end{equation}
This remains valid without any assertion about unrestricted social optimization. To compare with the complete game's $\OPT$, the following externality control is needed.

\begin{lemma}[Base social externality]\label{lem:externality}
If the base degree is at most $p$, then for a unilateral deviation $s\to t$ by a player of weight $w_i$,
\begin{equation}\label{eq:externality}
 W^B(t)-W^B(s)\le(p+1)w_i
 \sum_{e\in t_i\setminus s_i}c_e(x_e(t))
 \le(p+1)w_iC_i^B(t).
\end{equation}
\end{lemma}
\begin{proof}
Retained resources have the same load; deleted resources cannot increase $xc_e(x)$. At an added resource, with new load $y=x+w_i$, a degree-$k$ monomial contributes a social increment $y^{k+1}-x^{k+1}=w_i\sum_{r=0}^ky^{k-r}x^r\le(k+1)w_iy^k$. Sum with the nonnegative coefficients and then over added resources. This counts the externality on compulsory users as part of the true social cost.
\end{proof}
\begin{theorem}[Socially decreasing repair]\label{thm:socialrepair}
If every class with $n_a\ge3$ satisfies
\begin{equation}\label{eq:SD}
 3v_a-2u_a-z_a>(p+1)B_a,
\end{equation}
then every specified earliest-error correction strictly decreases $W^+$. Every global social optimum lies in $\F$, and
\begin{equation}\label{eq:optcompression}
 \OPT(G_+)=K_0+\min_{s\in\F}W^B(s).
\end{equation}
If~\eqref{eq:RD} also holds and $s\in\mathcal T_\alpha$ satisfies
$W^B(s)\le\lambda\min_{t\in\F}W^B(t)$ for $\lambda\ge1$, then $s$ is an $\alpha$-$\PNE$ and
\begin{equation}\label{eq:qualitytransfer}
 W^+(s)\le\lambda\OPT(G_+).
\end{equation}
An additive base-core gap $\varepsilon$ likewise transfers as an additive full-game gap $\varepsilon$.
\end{theorem}
\begin{proof}
The correction removes at least one trigger and creates none: the proof of Theorem~\ref{thm:core} shows every gate involving the mover becomes untriggered, and no other gate changes. The social gate reduction is at least $w_ad_a$. Lemma~\ref{lem:externality} bounds the base increase by $(p+1)w_aB_a$, so~\eqref{eq:SD} makes the total strictly decrease. No outside state can be optimal; on $\F$ use~\eqref{eq:allsocial} to obtain~\eqref{eq:optcompression}. For quality, stability follows from~\eqref{eq:toleranceset} and
\[
 W^+(s)=K_0+W^B(s)
 \le K_0+\lambda\min_{\F}W^B
 \le\lambda(K_0+\min_{\F}W^B),
\]
since $K_0\ge0$. The additive version uses the same equality. No division by the optimum is needed.
\end{proof}
When $3\alpha\ge p+1$, strong dominance~\eqref{eq:RD} already implies~\eqref{eq:SD}. Indeed $L_a\ge1$, $z_a\le u_a$, and
\[
 3v_a-2u_a-z_a>
 3\alpha B_a+(3\alpha L_a-3)u_a\ge(p+1)B_a.
\]
In particular every strongly dominating repair on a base of degree at most two has this social property. For higher-degree bases, an individually improving correction need not reduce social cost, which is why~\eqref{eq:SD} is retained. The degree in that condition is the base degree, even when the gate has higher degree.

\begin{corollary}[Regularized core selection with full-game quality]\label{cor:jointselection}
Suppose base resource functions $F_e$, with $F_e(\varnothing)=0$, have common insertion scales $h_i>0$ and factor $\alpha_0\ge1$:
\[
 h_ic_e(Y)\le\Delta_iF_e\le\alpha_0h_ic_e(Y)
\]
on their potential-user insertion domain, with actual load $Y$ after insertion. Write $\Phi_B=\sum_eF_e$, and suppose $0<m\le M<\infty$ satisfy $mW^B\le\Phi_B\le MW^B$ on $\F$. For $\eta\ge0$ put
\[
 \beta_\eta=\max_i\frac{\alpha_0h_i+\eta(p+1)w_i}{h_i+\eta w_i},\qquad
 \lambda_\eta=\frac{M+\eta}{m+\eta}.
\]
Choose a global minimizer $s_\eta$ of $\Phi_B+\eta W^B$ over $\F$. If the gates satisfy~\eqref{eq:RD} at $\alpha=\beta_\eta$ and satisfy~\eqref{eq:SD}, then $s_\eta$ is a full-game $\beta_\eta$-$\PNE$ and $W^+(s_\eta)\le\lambda_\eta\OPT(G_+)$.
\end{corollary}
\begin{proof}
The social resource function $xc_e(x)$ has insertion bounds $w_ic_e(Y)$ and $(p+1)w_ic_e(Y)$, as in the proof of Proposition~\ref{prop:regularization}. Consequently $\Phi_B+\eta W^B$ has common lower scale $h_i+\eta w_i$ and upper scale at most $\beta_\eta$ times it. Applying the deviation calculation of Lemma~\ref{lem:glue} to endpoints in $\F$ shows that a strict core $\beta_\eta$ improvement would decrease this objective. Its global core minimum is thus core stable and lifts by Theorem~\ref{thm:core}. If $o$ minimizes $W^B$ over $\F$, global minimality gives
\[
 (m+\eta)W^B(s_\eta)\le\Phi_B(s_\eta)+\eta W^B(s_\eta)
 \le\Phi_B(o)+\eta W^B(o)\le(M+\eta)W^B(o).
\]
Theorem~\ref{thm:socialrepair} transfers this quality ratio to the complete game.
\end{proof}
The required cost envelope concerns the real $W^B$. If optional-user potentials are normalized after absorbing compulsory loads, their constants must first be restored and that envelope proved again; a positive resource used only by compulsory players otherwise has zero compressed potential and positive social cost. The statement is a finite selection theorem. A core local minimum gives its stability conclusion, but the cost conclusion requires the stated global minimum or a separate global objective bound. In particular its gate dominance must be checked at $\beta_\eta$, rather than at an earlier factor before regularization.

\subsection{A positive selection example and a decision procedure}
Three unit-weight players choose resource $A$ of latency $x$ or resource $B$ of constant latency two; bit zero means $A$. Add the one anchored gate $c(x)=6(x/3)^2$ in the order $0,1,2$. Here $B_a=3$, $L_a=1$, $(z_a,u_a,v_a)=(2/3,8/3,6)$. At $\alpha=1$, the dominance margin is $6-3-8/3=1/3>0$; the social penalty is $18-16/3-2/3=12>(p+1)B_a=6$. The gate baseline is $K_0=6$.

The six front states have costs and exact stability as follows; the two excluded states are shown to make the full-game comparison explicit.
\begin{center}
\begin{tabular}{cccc}
\toprule
State & In $\F$ & $W^+$ & Exact equilibrium\\
\midrule
$000$ & yes & $15$ & no\\
$001$ & yes & $12$ & yes\\
$010$ & no & $24$ & no\\
$011$ & yes & $11$ & yes\\
$100$ & yes & $12$ & yes\\
$101$ & no & $23$ & no\\
$110$ & yes & $11$ & yes\\
$111$ & yes & $12$ & no\\
\bottomrule
\end{tabular}
\end{center}
These claims can also be checked without the table. Within the base core, a stable state has one or two users on $A$, with base total cost five or six respectively; all-$A$ and all-$B$ have preserving improvements. Theorems~\ref{thm:core} and~\ref{thm:socialrepair} then show that the full-game optimum is eleven and is achieved at an exact equilibrium. Selecting among the stable states reduces twelve to eleven while retaining exact stability. The output degree is two, whereas the base degree is one.

More generally, if the gates are already present in an instance and~\eqref{eq:RD} holds, one can enumerate $\F$, check the tolerance conditions~\eqref{eq:toleranceset}, and minimize the base cost on the resulting stable set, adding back $K_0$. With~\eqref{eq:SD}, the same domain also computes the true social optimum. The size of the resulting enumeration is $|\F|$, which need not be polynomial in a compact input. If a core descent potential is available, first use the specified repair and then follow strict full-game improvements inside $\F$. By~\eqref{eq:coremargin} every such core improvement is a strict base $\alpha$ improvement, so that potential proves termination. Without it, finite enumeration remains the justified method. The required existence and quality bounds on the front domain are supplied by the separate potential and optimization arguments.

\subsection{Interpreting the resulting certificates}
The hypotheses of the positive results describe different features of an instance.\label{sec:conclusion}
Small actual shares yield the integral-potential guarantee even without a restriction on the number of weight types. Finite marginal calibration instead uses the actual two-type populations and their scale constraints; player height offers a third bound for resources with many types. For complete anchored gates, the stable-set formula~\eqref{eq:toleranceset} and the social decomposition~\eqref{eq:allsocial} separate stability from the additional condition~\eqref{eq:SD} needed to compare with unrestricted social optimization.

In the near-degree construction, endpoint amplification, return windows, and same-degree gates have distinct jobs: they transmit a strict margin, cover the remaining front configurations, and exclude nonfront states. Their combination, rather than failure of a scale contract alone, gives the all-state lower bound.

\appendix
\section{Parameter estimates and complete state coverage}\label{app:construction}
This appendix verifies the uniform family defined in Section~\ref{sec:application}. We first bound the endpoint and return sums, then cover all relative phases of the front domain, and finally verify that the legal gates exclude every remaining state at the same degree. All costs include the mover's own weight.

\subsection{Endpoint, propagation, and return estimates}
We begin with the estimates used to transmit an improvement across adjacent weight classes and to close the return window.
\begin{lemma}[Endpoint and propagation estimates]\label{lem:appestimates}
The parameters satisfy $H,J>0$, $1<S\le bR<R$, and
\begin{equation}\label{eq:appJbounds}
 J/H\ge K/8>1,\qquad
 Y/J\le\theta/16,\qquad \alpha X\le2J,\qquad
 \kappa J-\alpha X\ge\theta J/2.
\end{equation}
Furthermore,
\begin{equation}\label{eq:appreturnbounds}
 g\le13n\log d,\qquad Z\ge N/2,\qquad
 \frac{A_RZ}{A_{g-1}J}\ge8,
\end{equation}
where $A_{g-1}=S^{g-1}$, and
\begin{equation}\label{eq:appdefectbounds}
 N\ge24(d+1)(gK+1)>gK,\qquad
 (\alpha+1)(t+1)W_0\le Z/4\quad(0\le t\le gK).
\end{equation}
\end{lemma}
\begin{proof}
Lemma~\ref{lem:area} gives
\begin{equation}\label{eq:appRiemann}
 R\frac{(1+\theta)T-1}{d+1}
 \le X\le R\frac{(1+\theta)T-1}{d+1}+T-1.
\end{equation}
The second condition in~\eqref{eq:uniformconditions} implies $\alpha/T\le\delta/16$, and thus $H\ge(1-\delta/16)T>T/2$. Also $d/R\le\delta/32$, and
\[
 (1-\delta)(1+\delta/8)^2\le1-3\delta/4,
 \qquad\kappa(1-\delta)<1.
\]
Using $J\le\alpha X$ in~\eqref{eq:appRiemann},
\begin{equation}\label{eq:appcontraction}
 \frac SR\le
 \frac{\kappa(1-\delta)(1+\theta+d/R)}{1-\alpha/T}
 \le\frac{1-23\delta/32}{1-\delta/16}
 \le1-\delta/2=b.
\end{equation}
Termwise comparison yields $Y\le\rho X$, with $\rho=(\theta/(1+\theta))^d\le\theta^2$. Since $\alpha\ge2$, $J\ge\alpha X/2>0$. The left side of~\eqref{eq:appRiemann} is at least $KT/(d+1)$, as $(1+\theta)T-1\ge\theta T$. Hence $J/H\ge\alpha K/(2(d+1))\ge K/8$, and $S>1$. Moreover
\[
 \frac YJ\le\frac{\theta^2}{2-\theta^2}\le\theta^2\le\theta/16,
\]
which proves all of~\eqref{eq:appJbounds} by $\alpha X=J+Y$ and $\kappa=1+\theta$.

Since $-\log b\ge1/(2n)$, minimality of $g$ gives
\[
 2\le g\le2n\log(192d^2)+1\le13n\log d.
\]
For the last inequality, $d\ge4$ suffices. The elementary exponential bound gives $W_0<3$. Every term in $Z$ is at least $1-\alpha d^{-d}\ge1/2$, so $Z\ge N/2$. As $H=\kappa J/S$,~\eqref{eq:appcontraction} implies
\[
 \frac{A_RZ}{A_{g-1}J}
 =\frac{\kappa Z}{4\alpha W_0S^g}
 \ge\frac1{24d^2b^g}\ge8.
\]
For $g\ge2$, $R^{g-2}\ge2^{g-2}\ge g/2$ by induction. Therefore
\[
 N\ge\frac{gR^2}{2d}\ge128gd^7
 \ge24(d+1)(gK+1)>gK.
\]
For example, the penultimate comparison follows from
$24(d+1)(gK+1)\le96gd^5$. Finally $\alpha\le d$ and $W_0<3$ give
\[
 (\alpha+1)(t+1)W_0\le3(d+1)(gK+1)\le N/8\le Z/4,
\]
proving~\eqref{eq:appdefectbounds}.
\end{proof}

\subsection{All front configurations have a preserving exit}\label{app:frontcoverage}
A front in class $a$ is represented by $x_a\in\mathbb Z/(2N)\mathbb Z$ through
\begin{equation}\label{eq:appfront}
 B_x(j)=\begin{cases}
 \ind[j<x],&0\le x<N,\\
 \ind[j\ge x-N],&N\le x<2N.
 \end{cases}
\end{equation}
Forward $x\to x+1$ flips clone $x\bmod N$; backward $x\to x-1$ flips clone $(x-1)\bmod N$. Thus these are genuine one-player deviations, not moves by a whole class.

We first justify a common phase normalization for the base window game from its real resource incidences. Shift all clone indices by one modulo $N$ and flip precisely those player bits crossing $N-1\to0$. If $e_h,e_u$ are the corresponding wrap indicators, the ordering predicate in~\eqref{eq:appincidence} changes by $e_h\xor e_u$. The light player's changed bit cancels $e_u$, and all users of an anchor pair change side by the same $e_h$. Exchange that pair's resource labels when $e_h=1$. The two sides have the same latency, so base-game loads and player costs, including costs after any specified single-player flip, are preserved. In~\eqref{eq:appfront} this transformation adds one to every phase. We may therefore normalize a chosen mover class to phase zero while leaving the relative phases arbitrary.

\begin{lemma}[Adjacent endpoint costs]\label{lem:appendpoints}
Normalize the heavy phase of an adjacent test to zero, and write $r=x_{a+1}-x_a\pmod{2N}$. Its heavy forward mover is clone zero, with old and new costs
\begin{equation}\label{eq:appheavy}
 A_a(1+h(r)/R)^d,\qquad
 A_a(1+(K-h(r))/R)^d,
\end{equation}
where
\[
 h(r)=\begin{cases}
 K-\min\{r,K\},&0\le r\le N,\\
 \min\{r-N,K\},&N\le r<2N.
 \end{cases}
\]
The heavy backward costs are the reverse of~\eqref{eq:appheavy} at $r+1\pmod{2N}$. If $r\notin[0,K]\cup[N,N+K]$, there is a heavy front-preserving direction with costs $(A_aT,A_a)$.

At any relative phase, either front-preserving endpoint flip of a light class has old and new costs at most $A_aX$ in its incoming adjacent test.
\end{lemma}
\begin{proof}
For heavy clone zero all its window users have greater global position. The negative orientation in~\eqref{eq:appincidence} therefore makes their resource sides equal their bits. Count the zero bits of $B_r$ on positions $0,\ldots,K-1$ to obtain $h(r)$ and the two complementary loads in~\eqref{eq:appheavy}. A backward edge is the reversal of the forward edge starting one phase earlier; relative to its normalized starting heavy phase, the light phase is $r+1$. For $K<r<N$ use backward, and for $N+K<r<2N$ use forward, obtaining the displayed costs.

For the light statement normalize its phase to zero and consider its forward mover, clone zero. It belongs to exactly $K$ windows, whose anchors have indices $-k\bmod N$, $0\le k<K$. Among the light users in each window, the number on its current side is $K-k$ and on its target side after moving is $k+1$, including the mover. For $k=0$ the anchor is clone zero and all window positions precede a wrap; for $k>0$ the window wraps through clone zero, with its $k$ high-index light users on the other side. This also gives the same counts at the wrap. Let $e_k\in\{0,1\}$ say whether that window's heavy anchor is on its current side. Its two costs are precisely
\begin{equation}\label{eq:applight}
 A_a\sum_{k=0}^{K-1}((K-k)/R+e_k)^d,\qquad
 A_a\sum_{k=0}^{K-1}((k+1)/R+1-e_k)^d.
\end{equation}
Both are at most $A_aX$. A backward edge reverses a forward edge at another relative phase, so the same bounds hold.
\end{proof}

If an adjacent difference is outside the two short intervals, choose the heavy exit from Lemma~\ref{lem:appendpoints}. For class zero the only other test is its one return pair; its new normalized load is at most $1+N\varepsilon$, so discarding nonnegative old return costs leaves total margin at least
\begin{equation}\label{eq:approotexit}
 H-\alpha A_RW_0=3H/4.
\end{equation}
For a heavy mover in class $1\le a<g$, the only other test is incoming. The corresponding total margin is at least
\begin{equation}\label{eq:appmiddleexit}
 A_aH-\alpha A_{a-1}X
 =A_{a-1}(\kappa J-\alpha X)
 \ge(\theta/2)A_{a-1}J\ge\theta H/2.
\end{equation}
The last step uses $S>1$ and $J>H$.

It remains to cover every front vector whose adjacent differences all lie in the short intervals. Write
\begin{equation}\label{eq:appshortdifferences}
 x_{a+1}-x_a=k_aN+t_a\pmod{2N},\quad
 k_a\in\{0,1\},\quad0\le t_a\le K.
\end{equation}
Summing gives $x_g-x_0=kN+t\pmod{2N}$, with a parity $k\in\{0,1\}$ and $0\le t=\sum_at_a\le gK<N$.

Normalize the last class to phase zero and write $r=x_0-x_g$. For its forward clone zero in the return test, index all $N$ anchors by $h_j=-j\bmod N$. Let $\sigma_0=1$, $\sigma_j=0$ for $j>0$, and $e_j(r)=\ind[B_r(h_j)=\sigma_j]$. By the positive orientation in~\eqref{eq:appincidence}, the unscaled old and new return costs are
\begin{equation}\label{eq:appreturncosts}
 O(r)=\sum_{j=0}^{N-1}((N-j)\varepsilon+e_j(r))^d,\qquad
 Q(r)=\sum_{j=0}^{N-1}((j+1)\varepsilon+1-e_j(r))^d.
\end{equation}
Indeed the light-side counts follow the same wrap count as~\eqref{eq:applight}, now with all $N$ light users. The all-one $e$ vector gives $(U,V)$ and the all-zero vector gives $(V,U)$. Flipping a single entry changes each cost by at most $W_0$, because every load is at most $1+N\varepsilon$.

If $k=0$, choose forward, with $r=-t\pmod{2N}$. At $t=0$ only $e_0$ is zero; at $1\le t<N$, precisely $e_0,\ldots,e_t$ are zero. Hence the vector differs from all ones in at most $t+1$ entries. If $k=1$, choose backward, with $r=N-t$. Its old/new costs are $Q(r+1),O(r+1)$. At $t=0$ the vector for $r+1=N+1$ is all zero; at $1\le t<N$ it has exactly $t$ ones. Thus in this direction it differs from all zeros in at most $t+1$ entries. Both parity cases yield a return margin at least
\begin{equation}\label{eq:appreturnmargin}
 A_R[Z-(\alpha+1)(t+1)W_0]\ge3A_RZ/4.
\end{equation}
The mover's other cost is the incoming adjacent test, whose new cost is at most $A_{g-1}X$. Lemma~\ref{lem:appestimates} therefore gives total margin at least
\begin{equation}\label{eq:appfinalfrontmargin}
 3A_RZ/4-\alpha A_{g-1}X
 \ge6A_{g-1}J-2A_{g-1}J
 =4A_{g-1}J>\theta H/2.
\end{equation}
The mismatch case and the short-difference case exhaust all phase vectors. We have proved that every product-front state has a preserving real single-player exit with margin at least
\begin{equation}\label{eq:appeta}
 \eta=\theta H/4>0.
\end{equation}

\subsection{All-state cost cap and same-degree completion}\label{app:samedegree}
At an arbitrary base state, each adjacent resource has potential load at most $(1+\theta)w_a$. A heavy user uses one pair and a light user uses $K$ pairs. Thus the contribution from any such test is at most $KA_aT$. With $1<S<R$, all adjacent tests together cost any player at most $gKS^{g-1}T\le gR^gT$. The return contribution is at most $NA_RW_0=NH/(4\alpha)\le R^gT$. Consequently
\begin{equation}\label{eq:appbasecap}
 C_i^B(s)\le B_*=(g+1)R^gT
 \quad\text{for all players and all base states},\qquad
 B_*/\eta\le64n(g+1)R^g.
\end{equation}
The complete anchored gates were specified after~\eqref{eq:maincompletion}. Their number per class is
\begin{equation}\label{eq:appclocks}
 L_c=\binom{N-1}{2}.
\end{equation}
To apply the safe budget~\eqref{eq:safebudget} at $D=d$, it remains to prove
\begin{equation}\label{eq:apprepairtarget}
 (2/3)^d\le\eta/(4\alpha^2B_*L_c).
\end{equation}
Since $L_c\le N^2\le R^{2g}$,~\eqref{eq:appbasecap} yields
\begin{equation}\label{eq:appGamma}
 4\alpha^2B_*L_c/\eta\le
 \Gamma_*=256nd^2(g+1)R^{3g}.
\end{equation}
The first uniform condition and $d\ge16$ imply
\[
 \log R\le5\log d,\qquad g\le13n\log d,\qquad
 g+1\le14n\log d\le14d\log d\le d^3.
\]
It follows that
\begin{equation}\label{eq:appGammalog}
 \log\Gamma_*\le195n(\log d)^2+8\log d
 \le200n(\log d)^2<d\log(3/2).
\end{equation}
Here $\log(256n)+5\log d\le8\log d$ follows from $n\le d/8$ and $d\ge16$, and $8\log d\le5n(\log d)^2$. The last strict inequality is exactly the third condition in~\eqref{eq:uniformconditions}. Thus~\eqref{eq:apprepairtarget} holds.

For clarity, this budget covers both kinds of state directly. On a front, the preserving exit of margin $\eta$ acquires new gate cost at most $L_cM(2/3)^d$, and its full margin remains at least $\eta/2$. On a nonfront, flip the earliest nonanchor bit different from the anchor back to that anchor. It leaves a triggered gate of cost $M$, all its new gates are untriggered, and its new base cost is at most $B_*$. Its full margin is at least
\[
 M-\alpha B_*-\alpha L_cM(2/3)^d
 \ge\alpha B_*-\eta/2\ge\eta/2>0.
\]
This is one deviation in the current state. The last inequality uses $B_*\ge\eta$. Every state therefore has a strict $\alpha$ improvement. The player set and action labels are unchanged, every new coefficient is positive rational, and the resource count is
\begin{equation}\label{eq:appresourcecount}
 2(g+1)N+2(g+1)\binom{N-1}{2}.
\end{equation}
Every old and new latency has degree exactly $d$. This proves Theorem~\ref{thm:uniformfamily} with a finite exact construction; the output can be very large.

\subsection{A uniform substitution approaching the degree}\label{app:tail}
We now check the parameter conditions when the approximation factor approaches $d$ with an explicit additive error. For an arbitrary integer $d\ge2^{20}$, put
\begin{equation}\label{eq:appnearparams}
 \ell=\lceil\log_2d\rceil,\qquad
 n=\left\lceil\frac{d}{1200\ell^2}\right\rceil.
\end{equation}
We check all three uniform conditions, including changes of $\ell$ at powers of two.
First $d>1200\ell^2$ throughout this range. At $d=2^{20}$ this is $2^{20}>1200\cdot20^2$. For $\ell\ge21$, use $d>2^{\ell-1}$ and $2^{20}>1200\cdot21^2$; the latter comparison propagates because $2(\ell/(\ell+1))^2\ge1$ for $\ell\ge3$. Therefore $n\ge2$ and
\begin{equation}\label{eq:approunding}
 \frac{d}{1200\ell^2}\le n
 \le\frac{d}{1200\ell^2}+1
 \le\frac{d}{600\ell^2}.
\end{equation}
In particular $d\ge8n$. The integral inequality $\log(1+x)\ge x/(1+x)$ gives
\[
 \log T\ge\frac{d}{8n+1}\ge\frac{d}{9n}
 \ge\frac{200}{3}\ell^2.
\]
But $n\le d$, $\log d\le\ell\log2<\ell$, and $\log16<4$, so
\[
 \log(16nd)\le\log16+2\log d<4+2\ell
 <\frac{200}{3}\ell^2\le\log T.
\]
Hence $T>16nd$. Finally~\eqref{eq:approunding} gives
\[
 200n(\log d)^2
 \le\frac d3(\log d/\ell)^2
 \le\frac d3(\log2)^2
 <d/3<d\log(3/2).
\]
The last inequality follows by integrating $1/t$ on $[1,3/2]$, giving $\log(3/2)>1/3$. All of~\eqref{eq:uniformconditions} holds. The uniform family has target
\[
 \alpha=d-d/n\ge d-1200\ell^2.
\]
Since all deviation costs are nonnegative, a strict improvement at factor $\alpha$ is strict at the smaller factor $d-1200\ell^2$ as well. This proves Theorem~\ref{thm:near} for every stated integer degree.

\subsection{Fixed linear targets and depth accounting}\label{app:fixedc}
When $n$ is fixed, the same estimates also give an explicit degree threshold and the exact number of weight levels.
\begin{corollary}\label{cor:fixedn}
For each fixed integer $n\ge2$, the construction works for every integer
\[
 d\ge d_0(n)=\max\{2^{20},(600n)^2\}.
\]
It has exactly $g+1\le13n\log d+1$ movable weight values and potential decreasing-weight height $g+1$.
\end{corollary}
\begin{proof}
These $d$ satisfy $d\ge8n$ and $d\ge512n^2$. A cubic binomial term gives
\[
 T\ge\binom d3(8n)^{-3}
 \ge\frac{d^3}{12288n^3}\ge16nd;
\]
the last inequality follows from $d^2\ge196608n^4$. For $d\ge2^{20}$, $\log d\le d^{1/4}$: it holds at the endpoint and $d^{1/4}-\log d$ increases for $d\ge256$. Hence
\[
 d\log(3/2)>d/3\ge200n\sqrt d
 \ge200n(\log d)^2.
\]
Theorem~\ref{thm:uniformfamily} applies. There is one weight $R^{-a}$ per class and the gates add no players. The adjacent resource pair anchored at $P_{a,0}$ has $P_{a+1,0}$ among its optional users, so $P_{0,0},\ldots,P_{g,0}$ is a decreasing-weight path of length $g+1$ in vertices. No longer decreasing path exists because there are only $g+1$ weights. This proves the exact count and height.
\end{proof}

\bibliographystyle{ACM-Reference-Format}
\bibliography{structural_methods_article}
\end{document}
